\ssscp{Distribution of apparent reflexes of **monsoŋ ‘green-blue’}{15-04-02-01}
Another example that straddles the Blust line that 
shows similar properties to **mantəR ‘cuscus’
is **monsoŋ `green-blue', seen in \trf{15-01}.

\begin{table}[p]\centering\tcp{Apparent reflexes of **monsoŋ ‘green-blue’}{15-01}
\begin{adjustbox}{totalheight=\textheight}
	\begin{threeparttable}\st{4pt}
	\begin{tabular}{l|llll}\lsptoprule
Language	&	Form	&	Gloss	&	Node	&	Comment\\
	&	**monsoŋ	&	green-blue	&		&	\\	\midrule
\ili{Makasar}	&	{\hr\hr}monʧoŋ	&	mountain	&	\tsc{SSul}	&	\cite[116]{GrimesGrimes1987}\\
\ili{Makasar}	&	{\hr\hr}monʧoŋ bulo	&	green	&	\tsc{SSul}	&	lit: ‘mountain fur’;\\
				&											&				&								&	\cite[126]{GrimesGrimes1987}\\
Konjo	&	{\hr\hr}monʧoŋ	&	green	&	\tsc{SSul}	&	\cite[126]{GrimesGrimes1987}\\
\ili{Selayar}	&	{\hr\hr}monʧoŋ	&	green	&	\tsc{SSul}	&	\cite[126]{GrimesGrimes1987}\\	\midrule
\ili{Laiyolo}	&	{\hr\hr}monsoŋ	&	green	&	\tsc{Cel}	&	\cite[127]{GrimesGrimes1987}\\	\midrule
\ili{Bima}	&	{\hr\hr}moro	&	greenish	&	\tsc{BL}	&	\tsc{East Sumbawa}\\
\ili{Kambera}	&	{\hr\hr}muru	&	green, fresh	&	\tsc{BL}	&	\tsc{Sumba}; \cite[295]{Onvlee1984}\\
\ili{Anakalang}	&	{\hr\hr}moru	&	green	&	\tsc{BL}	&	\tsc{Sumba}\\
\ili{Mamboru}	&	{\hr\hr}moru	&	green	&	\tsc{BL}	&	\tsc{Sumba}\\
\ili{Kodi}	&	{\hr\hr}moro	&	green	&	\tsc{BL}	&	\tsc{Sumba}\\	\midrule
\ili{Rikou}	&	{\hr\hr}modo-ʔ	&	green	&	\tsc{TB}\su{†}	&	\tsc{Nuclear Rote}\\
\ili{Lole}	&	{\hr\hr}mo{\tl}modo-k	&	green	&	\tsc{TB}	&	\tsc{Nuclear Rote}\\
\ili{Termanu}	&	{\hr\hr}mo{\tl}modo-k	&	yellow	&	\tsc{TB}	&	\tsc{Nuclear Rote}\\
\ili{Dela}	&	{\hr\hr}moɗo-ʔ	&	green, blue	&	\tsc{TB}	&	\tsc{West \tsc{Rote-Meto}}\\
\ili{Kotos} 	&	{\hr\hr}moro-ʔ	&	yellow	&	\tsc{TB}	&	\tsc{Meto}\\
\ili{Amarasi}	&		&		&		&	\\
\ili{Molo}	&	{\hr\hr}molo	&	yellow	&	\tsc{TB}	&	\tsc{Meto}\\
\ili{Tetun}	&	{\hr\hr}modo-k	&	yellow, green	&	\tsc{TB}	&	\tsc{Eastern Timor}\\
\ili{Galolen}	&	{\hr\hr}moso	&	green-blue	&	\tsc{TB}	&	\tsc{Atauro-Wetar}\\
\ili{Raklungu}	&	{\hr\hr}kah-moro	&	green-blue	&	\tsc{TB}	&	\tsc{Atauro-Wetar}\\
\ili{Perai}	&	{\hr\hr}mosoŋ	&	green, blue	&	\tsc{TB}	&	\tsc{Atauro-Wetar}\\
Ili{\Q}uun	&	{\hr\hr}mosoŋ	&	green, blue	&	\tsc{TB}	&	\tsc{Atauro-Wetar}\\
\ili{Tugun}	&	{\hr\hr}moso	&	green	&	\tsc{TB}	&	\tsc{Atauro-Wetar}\\
W. \ili{Damer}	&	{\hr\hr}meh{\tl}meh-o	&	green	&	\tsc{TB}	&	\tsc{SW. Maluku}\\
\ili{Emplawas}	&	{\hr\hr}mo{\tl}mo-i	&	green	&	\tsc{TB}	&	\tsc{SW. Maluku}\\
E. \ili{Damer}	&	{\hr\hr}mor{\tl}mor	&	green	&	\tsc{TB}	&	\tsc{SW. Maluku}\\
\ili{Roma}	&	{\hr\hr}mo{\tl}moti	&	green	&	\tsc{TB}	&	\tsc{SW. Maluku}\\
\ili{Kisar}	&	{\hr\hr}mok{\tl}moke	&	green	&	\tsc{TB}	&	\tsc{SW. Maluku}; **t > \it{k}\\
\ili{Leti}	&	{\hr\hr}mota	&	green	&	\tsc{TB}	&	\tsc{SW. Maluku}\\
\ili{Wetan}	&	{\hr\hr}mota	&	green	&	\tsc{TB}	&	\tsc{SW. Maluku}\\
\ili{Luang}	&	{\hr\hr}motə	&	green	&	\tsc{TB}	&	\tsc{SW. Maluku}\\
\ili{Makatian}	&	{\hr\hr}mos{\tl}mos-e	&	green	&	\tsc{TB}	&	\tsc{SW. Tanimbar}\\
S. \ili{Mambae}	&	{\hr\hr}moro	&	green	&	\tsc{CT}	&	\\
\ili{Kemak}	&	{\hr\hr}mosok	&	green-blue	&	\tsc{CT}	&	\\
		\lspbottomrule
	\end{tabular}
		\begin{tablenotes}\tnsize
			\item[†] {The \tsc{TB} words would be regular from intermediate **mozo.}
		\end{tablenotes}
	\end{threeparttable}
\end{adjustbox}
\end{table}

This etymon clearly straddles the Blust line; it has limited
distribution; it is found in some Wallacean subgroups
and not in others; it appears to be found in some nodes within
subgroups and not in others.
Its distribution strongly implies contact,
yet in many languages it appears to reflect the same or nearly
the same regular sound correspondences as one would expect from a PMP etymon.
Given that borrowed words are often based on apparent similarities
and do not always follow regular sound changes,
is this a loan from \ili{Makasar} sailors who have been trading across
the region for centuries? Is it influence from the Sultanate
of Goa (\ili{Makasar}) over \ili{Bima}
and beyond? 
The sporadic \tsc{\ili{Bima}-Lembata} and \tsc{Timor-Babar} data possibly
point to an earlier form **mozo (or possibly **mozaw).
But such forms are not compatible with \it{monʧoŋ/monsoŋ} in
the \ili{Makasar} languages since *-z- never becomes \it{-nʧ-},
and final \it{-oŋ} also does not match the \tsc{\ili{Bima}-Lembata}
and \tsc{Timor-Babar} data (except for \it{mosoŋ} in \ili{Perai}
and Ili{\Q}uun -- both of which favour a loan source).
However, there is a better match for **mozo in Proto-\ili{Bungku}-\ili{Tolaki}
*uso ‘green’ \citep[486]{Mead1998} or *moʔuso.
PMP *z > Proto-\ili{Bungku}-\ili{Tolaki} **s is regular,
as is PMP *-aw > *o.
These forms suggest an earlier reconstruction of *ma-quzaw ‘green-blue’
(which appears to be related to PMP *qizaw ‘green’).\footnote{
		While \citet{BlustTrussel2020} has PWMP *hizaw glossed as ‘fighting
		cock with greenish feathers on light background’,
		the glosses of the specific languages
		and the notes for this etymon in ACD allow ‘green-blue’.
		The additional data in the ACD notes also seem to imply wide
		discrepancies in their sound correspondences.
		\citet[982]{Wolff2010} makes this issue more explicit in his
		reconstruction of ††sijaw (= PMP *hizaw in the more general transcription)
		giving it a more generic gloss ‘green’
		and a note “Most attestations are irregular in their correspondences.
		This form has spread secondarily.” In other words,
		contact is a clear factor in the distribution
		and forms of these related etyma.
		We are grateful to an anonymous reviewer who pointed out further
		implications of the Sulawesi data
		and the Proto-\ili{Bungku}-\ili{Tolaki} forms to us.}

Regardless,
the form and distribution of this etymon in the western parts
of our target region are not greatly different from that of **mantəR
‘cuscus’ there, yet one would be hard-pressed to make convincing arguments for
this etymon to be the basis for a subgroup.
Given the diversity of the medial consonant in the daughter forms,
it is entirely possible that this was borrowed separately into several different subgroups.
Are we to posit a subgroup that includes \tsc{South Sulawesi},
\tsc{Celebic}, \tsc{\ili{Bima}-Lembata}, \tsc{Timor-Babar} and \tsc{\ili{Central}
Timor}? And perhaps posit a proto vowel *o as well? Again,
this illustrates the dangers of basing subgroups on lexical data
in a region of great mobility that does not dovetail with other stronger phonological evidence.

